\chapter{单侧极限、无穷极限与无穷远处极限}

\section*{A组　定义与单侧极限}
\addcontentsline{toc}{section}{A组　定义与单侧极限}

\begin{solution}{16.1}
右极限 \(f(x)\to L\) 的定义是
\[
 \forall\varepsilon>0\ \exists\delta>0\ \forall x\in D,\quad
 a<x<a+\delta\Rightarrow\abs{f(x)-L}<\varepsilon.
\]
左无穷极限 \(f(x)\to+\infty\) 是
\[
 \forall M>0\ \exists\delta>0\ \forall x\in D,\quad
 a-\delta<x<a\Rightarrow f(x)>M.
\]
在 \(+\infty\) 处趋于有限 \(L\) 是
\[
 \forall\varepsilon>0\ \exists A\in\R\ \forall x\in D,\quad
 x>A\Rightarrow\abs{f(x)-L}<\varepsilon.
\]
\end{solution}

\begin{solution}{16.2}
双侧极限推出两个单侧极限，只需限制量词范围。反向分别取得对同一 \(\varepsilon\) 有效的半径 \(\delta_-,\delta_+\)，取其最小值。删心邻域中的任意点必在左侧或右侧，故相应估计成立。两侧聚点保证左右极限都不是空泛条件，也使“存在且相等”的表述有实际内容。
\end{solution}

\begin{solution}{16.3}
若 \(x\to m-\)，充分靠近时 \(m-1<x<m\)，故 \(\lfloor x\rfloor=m-1\)。若 \(x\to m+\)，充分靠近时 \(m<x<m+1\)，故 \(\lfloor x\rfloor=m\)。因此
\[
 \lim_{x\to m-}\lfloor x\rfloor=m-1,\qquad
 \lim_{x\to m+}\lfloor x\rfloor=m.
\]
\end{solution}

\begin{solution}{16.4}
当 \(x<0\) 时 \(\abs x/x=-1\)，当 \(x>0\) 时为 \(1\)。故左右极限分别为 \(-1,1\)。两者不相等，所以双侧极限不存在。
\end{solution}

\begin{solution}{16.5}
给定 \(M>0\)，取 \(\delta=1/M\)。若 \(0<x<\delta\)，则 \(x<1/M\)，从而 \(1/x>M\)。这正是右侧趋于 \(+\infty\) 的定义。
\end{solution}

\begin{solution}{16.6}
给定 \(M>0\)，要求 \(3x+1<-M\)，等价于
\[
 x<-\frac{M+1}{3}.
\]
取 \(A=-(M+1)/3\)，则所有 \(x<A\) 都满足所需不等式，故极限为 \(-\infty\)。
\end{solution}

\begin{solution}{16.7}
若定义对所有 \(M>0\) 成立，则给定 \(A\in\R\)，在 \(A>0\) 时取 \(M=A\)，在 \(A\le0\) 时取任意正阈值如 \(M=1\)，均得最终 \(f(x)>A\)。反之，“对每个实阈值”显然包含全部正阈值，故恢复原定义。
\end{solution}

\begin{solution}{16.8}
若 \(f(x)\to+\infty\)，则给定 \(M>0\)，最终有 \(f(x)>M\)，因而 \(\abs{f(x)}>M\)。反向不成立，例如 \(f(x)=1/x\) 在 \(x\to0\) 时满足 \(\abs{f(x)}\to+\infty\)，但左右符号相反，原函数没有双侧扩充实极限。
\end{solution}

\section*{B组　序列刻画、运算与代换}
\addcontentsline{toc}{section}{B组　序列刻画、运算与代换}

\begin{solution}{16.9}
量词否定为
\[
 \exists M_0>0\ \forall\delta>0\ \exists x\in D,\quad
 0<\abs{x-a}<\delta,\quad f(x)\le M_0.
\]
按 \(\delta=1/n\) 选点，得到 \(x_n\to a\) 且 \(f(x_n)\le M_0\) 对所有 \(n\) 成立，所以像列不趋于 \(+\infty\)。反之，存在这样的见证列会与扩展 Heine 定理冲突。
\end{solution}

\begin{solution}{16.10}
给定 \(M>0\)，取 \(\delta=1/\sqrt M\)，则 \(0<\abs x<\delta\) 推出 \(1/x^2>M\)。对 \(1/x\)，沿 \(x_n=1/n\) 趋于 \(+\infty\)，沿 \(y_n=-1/n\) 趋于 \(-\infty\)，所以既无有限极限，也不趋于同一个无穷方向。
\end{solution}

\begin{solution}{16.11}
由 \(g(x)\to L\) 的局部有界性，可在某相应邻域内取 \(\abs{g(x)}\le K\)。给定 \(M>0\)，由 \(f(x)\to+\infty\) 使 \(f(x)>M+K\)。在共同邻域内
\[
 f(x)+g(x)\ge f(x)-\abs{g(x)}>M.
\]
\end{solution}

\begin{solution}{16.12}
由 \(g(x)\to M>0\)，最终有 \(g(x)>M/2\)。给定阈值 \(A>0\)，由 \(f(x)\to+\infty\) 使 \(f(x)>2A/M\)。共同邻域内
\[
 f(x)g(x)>\frac{2A}{M}\frac M2=A.
\]
\end{solution}

\begin{solution}{16.13}
在 \(x\to+\infty\) 时取
\[
 (f,g)=(x+1,x),\qquad (2x,x),\qquad
 (x+\sin x,x).
\]
每一对的两个函数都趋于 \(+\infty\)，而差依次为 \(1\)、\(x\to+\infty\)、\(\sin x\) 无极限。
\end{solution}

\begin{solution}{16.14}
在 \(x\to+\infty\) 时统一取 \(g(x)=x\)，并依次取
\[
 f_1(x)=x^{-2},\quad f_2(x)=x^{-1},\quad
 f_3(x)=x^{-1/2},\quad f_4(x)=\frac{\sin x}{x}.
\]
则 \(f_j\to0\)，而乘积依次为 \(1/x\to0\)、\(1\)、\(\sqrt x\to+\infty\)、\(\sin x\) 无极限。
\end{solution}

\begin{solution}{16.15}
必要性：给定 \(\varepsilon>0\)，由函数极限取阈值 \(A\)；由 \(x_n\to+\infty\)，充分大时 \(x_n>A\)，故像误差小于 \(\varepsilon\)。

充分性用反证。若函数极限不为 \(L\)，存在 \(\varepsilon_0>0\)，使对每个 \(A\) 都可选 \(x\in D\)、\(x>A\) 且 \(\abs{f(x)-L}\ge\varepsilon_0\)。令 \(A=n\) 逐项选择 \(x_n\)，则 \(x_n\to+\infty\)，但像列不趋于 \(L\)，矛盾。
\end{solution}

\begin{solution}{16.16}
设 \(F(x)\to L\)（\(x\to+\infty\)）。给定 \(\varepsilon>0\)，取 \(A>0\) 使 \(x>A\) 时误差小于 \(\varepsilon\)，再令 \(\delta=1/A\)。若 \(0<t<\delta\)，则 \(1/t>A\)，故 \(\abs{F(1/t)-L}<\varepsilon\)。

反之，给定 \(\varepsilon>0\)，由 \(F(1/t)\to L\)（\(t\to0+\)）取 \(\delta>0\)，令 \(A=1/\delta\)。若 \(x>A\)，则 \(0<1/x<\delta\)，故 \(\abs{F(x)-L}<\varepsilon\)。
\end{solution}

\section*{C组　统一形式、边界与项目}
\addcontentsline{toc}{section}{C组　统一形式、边界与项目}

\begin{solution}{16.17}
若 \(f(x)\to0\) 且最终为正，给定 \(M>0\)，把零极限精度取为 \(1/M\)，便有
\[
 0<f(x)<1/M,\qquad1/f(x)>M.
\]
反之，若 \(1/f(x)\to+\infty\)，给定 \(\varepsilon>0\)，取阈值 \(1/\varepsilon\)，则最终
\[
 0<f(x)<\varepsilon.
\]
\end{solution}

\begin{solution}{16.18}
因 \(x^2\to0+\)（\(x\to0\)），若 \(F(u)\to L\)（\(u\to0+\)），复合极限给出 \(F(x^2)\to L\)。反向也成立：给定任意 \(u_n\to0+\)，取 \(x_n=\sqrt{u_n}\to0+\)，由 \(F(x_n^2)\to L\) 得 \(F(u_n)\to L\)，再用 Heine 定理。因而这两个极限等价；它们都不涉及 \(F\) 在 \(0\) 左侧的行为。
\end{solution}

\begin{solution}{16.19}
对有限 \(\alpha\)，邻域基可取 \((\alpha-\delta,\alpha+\delta)\)；对 \(+\infty\) 可取 \((A,+\infty]\)；对 \(-\infty\) 可取 \([-\infty,A)\)。目标 \(\beta\) 也有同样三类邻域。任取三种输入中心与三种输出中心组合，统一条件都是：每个目标邻域 \(V\) 的原像在相应意义下最终包含输入中心的某个删心邻域。九种组合正好对应有限或无穷自变量与有限或无穷函数值。
\end{solution}

\begin{solution}{16.20}
必要性对三种输入中心都由“点列最终进入每个输入邻域”推出。充分性失败时固定一个目标坏邻域 \(V\)：有限 \(\alpha\) 时在 \(0<\abs{x-\alpha}<1/n\) 中选坏点；\(+\infty\) 时在 \(x>n\) 中选坏点；\(-\infty\) 时在 \(x<-n\) 中选坏点。所得点列趋于相应 \(\alpha\)，像列却从不进入 \(V\)，与序列条件矛盾。
\end{solution}

\begin{solution}{16.21}
典型转换包括
\[
 x\to a\Longleftrightarrow h=x-a\to0,\qquad
 x\to+\infty\Longleftrightarrow t=1/x\to0+,
\]
\[
 x\to-\infty\Longleftrightarrow t=1/x\to0-,\qquad
 x\to+\infty\Longleftrightarrow t=e^{-x}\to0+.
\]
平移与倒数在相应定义域上给出等价；\(u=x^2\) 把 \(x\to0\) 映到 \(u\to0+\)，只给右极限信息；一般复合代换先给单向蕴含，只有代换对目标邻域中的允许方向具有充分覆盖性时才能反推。
\end{solution}

\begin{solution}{16.22}
例如映射
\[
 \Phi(x)=\frac2\pi\arctan x
\]
把 \(\overline{\R}\) 序同构到闭区间 \([-1,1]\)，约定 \(\Phi(\pm\infty)=\pm1\)。可定义
\[
 d_*(x,y)=\abs{\Phi(x)-\Phi(y)}.
\]
该度量的收敛恰好是扩充实数的序拓扑收敛：有限点处等价于通常收敛，趋于 \(+\infty\) 等价于 \(\Phi(x)\to1\)，负无穷类似。
\end{solution}

\begin{solution}{16.23}
令
\[
 L=\sup\{f(x):x>A\},
\]
该上确界有限。给定 \(\varepsilon>0\)，按上确界定义存在 \(x_0>A\) 使 \(f(x_0)>L-\varepsilon\)。由单调递增性，所有 \(x>x_0\) 满足
\[
 L-\varepsilon<f(x)\le L.
\]
故 \(f(x)\to L\)（\(x\to+\infty\)）。完备性通过上确界的存在进入证明。
\end{solution}

\begin{solution}{16.24}
当 \(x\to+\infty\) 时，
\[
 (x+1)-x\to1,\quad 2x-x\to+\infty,\quad
 x-(2x)\to-\infty,
\]
而 \((x+\sin x)-x=\sin x\) 无极限。对 \(0\cdot\infty\)，取无穷因子 \(g=x\)，零因子分别取 \(x^{-2},x^{-1},x^{-1/2},\sin x/x\)，乘积依次趋于 \(0,1,+\infty\) 或无极限。这些结果说明未定式没有预先指定数值。
\end{solution}
